% NichidaiMasterExam_R7_1st_1
\begin{tcolorbox} %%R7_1st_1
	$\fbox{1}$ $A= \begin{pmatrix}1& 1& 1\\ 1& 0& 1\\ {-1}& {-1}& {-1}\end{pmatrix}$とし、$X$を$n$次複素正方行列で、ある非負整数$k$に対して$X^k$が零行列となるとする。
\begin{enumerate}
	\item $A$の固有多項式・固有値を求めよ。
	\item $X$の固有値をすべて求めよ。
	\item $X^m$が零行列となる最小の非負整数を$m$とすると、$m\le n$であることを示せ。
	\item $X$が対角化可能ならば、$X$は零行列であることを示せ。
\end{enumerate}
\end{tcolorbox}

\begin{souanswer} %%R7_1st_1_ans
	$\fbox{1}$
\begin{enumerate}
	\item 固有値を$\lambda$とする。
\begin{align*}
	\begin{vmatrix}A- \lambda E\end{vmatrix}
&= 
\begin{vmatrix}
	1- \lambda& 1& 1\\
	1& -\lambda& 1\\
	-1& -1& -1- \lambda
\end{vmatrix}\\
&= \lambda(1- \lambda)(1+ \lambda)- \lambda\\
&= -\lambda^3
\end{align*}
\begin{tabular}{lll}
	固有方程式& $\colon$& $\begin{vmatrix}A- \lambda E\end{vmatrix}= -\lambda^3$\\
	固有値& $\colon$& 0
\end{tabular}

	\item 固有値を$\mu$、対応する固有ベクトルを$\bm{v}\ (\bm{v}\ne \bm{0})$とする。固有方程式より
\begin{equation*}
	X\bm{v} = \mu \bm{v}
\end{equation*}
	この式の両辺に左から繰り返し$X$をかけ続ける。ここで任意の非負整数$k$に対して$X^k= O$であったから
\begin{equation*}
	X^k\bm{v}= \mu^k \bm{v}\iff Ov= \mu^k \bm{v}\iff \bm{0}= \mu^k \bm{v}
\end{equation*}
	ここで$\bm{v}= \bm{0}$であるから$\mu^k= 0$となるしかない。$\mu$を複素数の範囲で解くと$\mu= 0$
\end{enumerate}
\end{souanswer}